%% TVCG submission --- VerbalVis
%% Compile: pdflatex main && bibtex main && pdflatex main && pdflatex main

\documentclass[journal]{IEEEtran}

\usepackage{cite}
\usepackage{graphicx}
\graphicspath{{png/}}
\usepackage{capt-of}
\usepackage{amsmath,amssymb,amsfonts}
\usepackage{booktabs}
\usepackage{array}
\usepackage{tabularx}
\usepackage{url}
\usepackage{xcolor}
\usepackage{hyperref}
\hypersetup{colorlinks=true,linkcolor=blue,citecolor=blue,urlcolor=blue}

\begin{document}

\title{VerbalVis: Full-Duplex Conversational Visual Analytics for Analytical Intent Revision}

\author{Anonymous Author(s)%
\thanks{Manuscript submitted to IEEE Transactions on Visualization and Computer Graphics. Source code, prompts, formative-inquiry materials, and case-study traces will be made available upon publication.}}

\IEEEaftertitletext{%
\begin{center}
\includegraphics[width=\textwidth]{teaser-main.png}
\captionof{figure}{VerbalVis coordinates user speech, assistant speech, analytical tools, and a persistent dashboard during intent revision. In the running example, an analyst examines whether \texttt{office\_furniture} should receive priority improvement resources in RJ. While the assistant is explaining a preliminary risk-based recommendation, the analyst interrupts to require product revenue and order volume to be considered. VerbalVis stops the obsolete audio, rejects late events from the interrupted response, interprets the revised request against the retained RJ dashboard state, executes the accepted tool batch, and returns a new explanation grounded in the updated views.}
\label{fig:teaser}
\end{center}
}

\maketitle

\begin{abstract}
Exploratory data analysis is iterative: visual evidence can redirect an analyst's question, qualify a provisional explanation, or change the data scope under investigation. Speech is well suited to expressing such changes quickly and conversationally, while visualization provides a persistent workspace for comparing values, locating anomalies, and verifying spoken interpretations. Existing conversational visual analytics systems, however, usually organize refinement around completed turns. When a user interrupts an assistant that is still speaking, stopping audio alone is insufficient because the same model response may also contain analytical tool calls that can modify a persistent dashboard.

We present \textbf{VerbalVis}, a full-duplex conversational visual analytics system that coordinates user audio, assistant audio, schema-grounded analytical tools, and dashboard state. The current prototype uses response-level cancellation, a current-response admission check for function calls, a non-preemptive sequential tool boundary, and post-execution dashboard-state reinjection. If user speech supersedes an active response, VerbalVis immediately stops playback, sends \texttt{response.cancel}, and discards late audio, transcript, and completion events from the interrupted response. Function calls are accepted only from the current response after it completes normally. Once a tool batch begins, user audio is temporarily blocked and the batch is allowed to finish; the resulting tool outputs and latest dashboard state are then returned to the realtime model before a new spoken response is requested.

A four-participant formative inquiry informed requirements for context grounding, composable analytical actions, overlap handling, and cross-stream coordination. We use an RJ delivery-resource decision as a running example and present an SP weekly operational-risk case study to demonstrate how coordinated visual comparison supports hypothesis revision and scope refinement. The evidence is case-based rather than causal: it shows how the implemented mechanism behaves and what analytical work it supports, but it does not establish population-level usability or performance effects.
\end{abstract}

\begin{IEEEkeywords}
Visual analytics, conversational interfaces, speech interaction, full-duplex dialogue, analytical intent revision, exploratory data analysis, case study.
\end{IEEEkeywords}

\section{Introduction}
\label{sec:intro}

\IEEEPARstart{E}{xploratory} data analysis (EDA) is inherently iterative. Analysts move among observations, questions, tentative explanations, and new analytical actions as evidence emerges. Patterns and anomalies can expose limitations in the current direction and prompt analysts to reconsider what they are asking, what they believe, or which data are relevant \cite{russell1993cost,pirolli2005sensemaking,sacha2014knowledge,battle2019characterizing}. We use \emph{analytical intent revision} to denote an episode in which an analyst supersedes, qualifies, or materially changes an active analytical commitment.

Language interfaces and LLM-powered visual analytics systems increasingly support contextual queries, visualization generation, conversational refinement, and multi-step tool use \cite{gao2015datatone,setlur2016eviza,narechania2021nl4dv,dibia2023lida,wang2024dataformulator2,zhao2024lightva}. These capabilities make analytical operations easier to express, but refinement is usually organized around completed turns. Speech makes this constraint particularly visible because an assistant's answer unfolds over time. An analyst may realize that the analysis is misdirected well before the answer finishes, yet a turn-based interface requires waiting before the new direction can be expressed.

Voice and visualization serve complementary roles. Voice lowers the interaction cost of posing follow-up questions, adding constraints, and redirecting an analysis; it can make the exchange feel more like a discussion than a sequence of form submissions. This benefit should not be confused with lower model inference latency: the advantage is that the analyst can express a change immediately. Visualization, in contrast, externalizes the current analytical state. It supports visual comparison across metrics, persistent inspection of evidence, and detection of cross-metric anomalies that are difficult to retain from speech alone. For example, a category may have the highest product revenue, a comparatively low low-score ratio, and a relatively small order count. The combination is more informative than any metric in isolation because it may indicate high unit value rather than broad sales volume.

Full-duplex speech allows the system to keep listening while speaking, so a user can redirect an ongoing explanation through barge-in. Yet barge-in alone does not solve the visual analytics problem. It is a temporal event, not necessarily evidence of analytical revision: users may interrupt to acknowledge the assistant, repair recognition, request clarification, or stop an overlong answer. More importantly, a realtime model response may contain both spoken output and one or more function calls. If the response becomes obsolete, the system must decide whether its speech should stop, whether its function calls should still be admitted, and which dashboard state should ground the next response.

We present \textbf{VerbalVis}, a full-duplex conversational visual analytics system designed around this coordination problem. VerbalVis does not introduce a new speech model. Speech understanding, server-side voice activity detection, speech generation, and function-call planning are provided by Qwen Realtime. The system contribution lies in what happens after provider events arrive: playback management, response validity, tool admission, dashboard updates, and context reinjection.

The implemented coordination policy has four parts. First, the system maintains exactly one active assistant response. On a server-side speech-start event, it stops browser playback and sends \texttt{response.cancel}. Second, late audio, transcripts, and completion events are accepted only when their response identifier still matches the current response. Third, function calls are extracted only from the current response after it completes normally; a cancelled or superseded response therefore cannot submit unfinished calls. Fourth, once a valid tool batch starts, it becomes an input-closed, non-preemptive boundary: microphone chunks are blocked, tools run sequentially, and their outputs plus the latest dashboard state are returned to the model before one follow-up response is requested. Running tools are not cancelled or rolled back.

Drawing on sensemaking and information-seeking research, we describe revisions through three theory-informed, non-exhaustive, and potentially overlapping objects of change: the analytical goal, a working hypothesis, and the analytical scope \cite{pirolli2005sensemaking,klein2007dataframe,bates1989berrypicking,battle2019characterizing}. These constructs support design reasoning and post-hoc interpretation; they are not runtime intent classes.

We retain a four-participant formative inquiry as design-oriented evidence. We then use two complementary analytic scenarios. The RJ delivery-resource decision is the \emph{running example} used throughout the design discussion. The SP weekly operational-risk analysis is a separate \emph{case study} used to demonstrate the complete system loop, cross-metric visual reasoning, and response-level redirection. We do not conduct a comparative user study and therefore make no causal claim that full-duplex interaction improves task performance, workload, or perceived control.

\textbf{Contributions.}
\begin{enumerate}
\setlength{\itemsep}{1pt}
\setlength{\parskip}{0pt}

\item A design framing for analytical intent revision in full-duplex conversational visual analytics, distinguishing speech overlap from changes to analytical goals, working hypotheses, and analytical scope.

\item \textbf{VerbalVis}, a full-duplex visual analytics system that combines response-level cancellation, current-response tool admission, non-preemptive sequential tool execution, and post-tool dashboard-state reinjection.

\item A case-based account of the system, including an RJ running example, an SP operational-risk case study, and mechanism-oriented checks for audio replacement, tool-flow consistency, dashboard synchronization, and metric semantics.
\end{enumerate}

\section{Related Work}
\label{sec:related}

VerbalVis connects research on natural-language interaction for visualization, LLM-powered visual analytics, full-duplex spoken dialogue, and analytical sensemaking.

\subsection{Natural-Language Interaction for Visualization}
\label{sec:rw-nli}

Natural-language interfaces for visualization translate utterances into data attributes, analytical tasks, and visual specifications. DataTone exposes alternative interpretations through interactive ambiguity widgets \cite{gao2015datatone}, while Eviza grounds language in both the dataset and active visualization to support contextual references and follow-up questions \cite{setlur2016eviza}. NL4DV provides a reusable pipeline for extracting attributes and analytical tasks and generating candidate visualizations \cite{narechania2021nl4dv}. These systems establish language grounding, ambiguity management, and contextual interpretation as central requirements.

Multimodal systems combine language with visible references and direct manipulation. Orko integrates speech and touch for network exploration \cite{srinivasan2018orko}; InChorus characterizes operations expressed through speech, pen, and touch \cite{srinivasan2020inchorus}; and DataBreeze interweaves speech with direct manipulation over flexible unit visualizations \cite{srinivasan2021databreeze}. Snowy and conversational extensions of NL4DV further support follow-up utterances and modification of existing views \cite{srinivasan2021snowy,mitra2022conversational}. These systems mainly organize refinement after an interaction turn completes, whereas VerbalVis focuses on redirection while assistant output is still unfolding.

\subsection{LLM-Powered and Agentic Visual Analytics}
\label{sec:rw-agentic}

Large language models extend conversational visualization beyond mapping one utterance to one chart. LIDA decomposes visualization generation into data summarization, goal generation, chart creation, and evaluation \cite{dibia2023lida}. Data Formulator and Data Formulator~2 combine language with explicit encodings and iterative histories \cite{wang2024dataformulator,wang2024dataformulator2}. InsightPilot and LightVA translate high-level goals into sequences of analytical tasks and visualizations \cite{ma2023insightpilot,zhao2024lightva}; WaitGPT externalizes intermediate operations \cite{xie2024waitgpt}; and InsightLens links conversation, code, visualizations, and findings across an extended analysis \cite{weng2025insightlens}.

These systems make analytical processes more steerable and inspectable, but steering is usually performed through later prompts, direct-manipulation controls, or explicit checkpoints. Existing work rarely specifies what happens when a user changes direction before a model response and its associated analytical actions have completed.

\subsection{Full-Duplex Dialogue and Interruption}
\label{sec:rw-duplex}

Streaming speech systems incrementally process or generate audio, whereas full-duplex systems continue listening while speaking. They must handle background speech, backchannels, cooperative overlap, and attempts to take the conversational floor \cite{skantze2021turntaking}. Barge-in traditionally refers to detecting user speech during system output and stopping or modifying the active response \cite{heins1997bargein}. Recent systems including Duplex Conversation, SyncLLM, Moshi, and SALMONN-omni investigate concurrent speech streams, low-latency turn-taking, and context-sensitive interruption \cite{lin2022duplex,veluri2024syncllm,defossez2024moshi,yu2025salmonnomni}.

Interruption timing does not determine analytical meaning. A user may interrupt to acknowledge, repair recognition, ask for clarification, or revise the investigation. Moreover, stopping audio does not automatically determine the status of function calls or dashboard updates. VerbalVis therefore treats speech onset as a floor-management signal and uses response identity plus a tool-execution boundary to coordinate downstream analytical work.

\subsection{Sensemaking and Analytical Revision}
\label{sec:rw-sensemaking}

Sensemaking research characterizes analysis as an iterative process in which questions, explanations, and data focus evolve as evidence is encountered \cite{russell1993cost,pirolli2005sensemaking,sacha2014knowledge}. Information-seeking studies describe evolving queries and the refinement of broad goals into focused analytical subtasks \cite{bates1989berrypicking,vakkari2001changes,battle2019characterizing}. Data--Frame Theory and scientific reasoning show that provisional explanations may be questioned or replaced when observations conflict with expectations \cite{klein2007dataframe,klahr1988dual,chinn1993anomalous}. Information-foraging and visualization-task frameworks distinguish analytical purpose from the population, time, geography, variable, category, and granularity under investigation \cite{pirolli1999information,brehmer2013multilevel,wongsuphasawat2016voyager}.

These traditions support the distinction among analytical goal, working hypothesis, and analytical scope, but they do not establish a universal or mutually exclusive taxonomy. We therefore treat the three constructs as sensitizing concepts rather than runtime labels.

\section{Formative Inquiry and Design Requirements}
\label{sec:formative}

We conducted a formative inquiry with four participants (F01--F04) to understand how analysts revise an ongoing investigation during conversational visual analysis and what system behavior such revisions require. The inquiry was design-oriented rather than confirmatory: it refined operational definitions, identified compound and boundary cases, and informed VerbalVis requirements.

\subsection{Theory-Informed Revision Constructs}
\label{sec:formative-constructs}

We specified three theory-informed concepts before analyzing the sessions:

\begin{description}
\item[\textbf{Analytical Goal Shift.}]
The analyst materially redirects the primary analytical question or desired knowledge outcome.

\item[\textbf{Working-Hypothesis Revision.}]
The analyst rejects, replaces, or qualifies an active provisional explanation, interpretation, or decision-relevant proposition.

\item[\textbf{Analytical Scope Refinement.}]
The analyst changes constraints over the relevant population, time period, geography, variables, categories, granularity, or data subset while generally preserving the higher-level analytical objective.
\end{description}

The constructs describe different objects of change rather than mutually exclusive event classes. A single utterance may receive multiple labels. Hypothesis formation, method changes, representation changes, evidence requests, and conversational repair were retained as boundary codes.

\subsection{Procedure and Analysis}
\label{sec:formative-procedure}

Participants used an early VerbalVis prototype to explore the Olist Brazilian e-commerce dataset through speech. The initial dashboard presented monthly order volume, category-level sales, state-level order distribution, review-score distribution, and delivery-time statistics. The exploration was open-ended; participants were not introduced to the revision constructs and were not required to interrupt the assistant or follow a prescribed path.

We recorded timestamped user and assistant utterances together with available dashboard events. After exploration, participants completed a short questionnaire about when and why they changed direction, which state should be preserved or abandoned, and which operations they considered important.

Our analysis followed a hybrid deductive--inductive process. An LLM-assisted pass retrieved candidate episodes from the logs, including the utterance, preceding commitment, proposed labels, timing, and rationale. We then manually reviewed each candidate against the surrounding dialogue and dashboard context. An episode was retained only when an active analytical commitment was identifiable and the utterance materially changed, qualified, or superseded it.

The final coding retained 27 analytical revision episodes. Under a multi-label scheme, 12 involved an Analytical Goal Shift, five involved a Working-Hypothesis Revision, and 18 involved an Analytical Scope Refinement. Eight involved more than one construct; therefore, counts do not sum to 27.

\begin{table*}[t]
\centering
\caption{Representative formative episodes. Utterances are translated from Chinese.}
\label{tab:formative-episodes}
\small
\begin{tabularx}{\textwidth}{p{0.17\textwidth}p{0.22\textwidth}Xp{0.25\textwidth}}
\toprule
\textbf{Construct} & \textbf{Prior commitment} & \textbf{User utterance} & \textbf{Interpretation} \\
\midrule
\textbf{Analytical Goal Shift} & Explore monthly order trends. & ``Analyze customer experience---the low-rated orders, including their time, category, and delivery conditions.'' & The desired outcome changes from describing order volume to investigating low-rating customer experience. \\
\textbf{Working-Hypothesis Revision} & Prioritize Bahia because it has the longest delivery time among high-volume states. & ``BA only has the longest delivery time. Does this take sales volume into account?'' & The user qualifies the decision criterion by requiring business volume to be considered. \\
\textbf{Analytical Scope Refinement} & Investigate customer experience over all orders. & ``Now filter the global review scores to one and two.'' & The objective is retained while the relevant population is restricted to low-rating orders. \\
\bottomrule
\end{tabularx}
\end{table*}

\subsection{Findings and Requirements}
\label{sec:formative-findings}

Participants revised different analytical commitments. Goal shifts changed the knowledge outcome; scope refinements changed the relevant population, time range, category, geography, or granularity; and working-hypothesis revisions questioned an active explanation or decision criterion. Eight episodes changed more than one analytical object, indicating that the constructs should not be implemented as mutually exclusive runtime states.

Speech overlap was also distinct from analytical revision. Participants spoke during output to correct recognition, ask follow-up questions, shorten an explanation, or acknowledge the assistant. Therefore, overlap indicates that the system should yield the floor, but the content of the completed utterance must determine the analytical action.

Participants expected selective cancellation and preservation. They wanted unfinished explanations and irrelevant pending work to be abandoned, while still-relevant filters, visualizations, findings, and dialogue context should remain available. These observations informed four requirements:

\paragraph{DR1: Ground interpretation in current dialogue and visual state.}
Follow-up requests should be interpreted relative to active filters, visible views, recent variables, and the aspect currently being explained.

\paragraph{DR2: Support compound revisions with composable analytical actions.}
One utterance may change scope, request a comparison, and create or modify a view. The system should permit multiple schema-grounded actions rather than force one exclusive intent label.

\paragraph{DR3: Separate floor yielding from semantic interpretation.}
Speech onset should stop the assistant from occupying the audio channel, but overlap alone should not be equated with an analytical revision.

\paragraph{DR4: Coordinate response validity, tool admission, and dashboard state.}
When an active response is superseded, obsolete speech should stop and unfinished function calls from that response should not be admitted. Already established dashboard state should be preserved unless the new request changes it. Once a valid tool batch starts, the prototype may complete it as a closed execution boundary and should ground the next response in the resulting dashboard state.

This inquiry provides design-oriented evidence, not validation of a universal taxonomy. Four participants cannot establish saturation or population-level prevalence. The subsequent running example and case study demonstrate the resulting system design rather than testing human-performance hypotheses.

\section{Running Example: RJ Delivery-Resource Allocation}
\label{sec:running-example}

We use an RJ resource-allocation scenario as a running example throughout the system description. The analytical scope is customer orders from RJ between October 2017 and May 2018. The business question is whether \texttt{office\_furniture} should receive priority delivery resources among the product-revenue Top~15 categories.

A preliminary risk-focused view may make \texttt{office\_furniture} appear compelling because it combines a high low-score ratio with long delivery time. The decision is not straightforward, however, because the category can also have relatively few orders. A manager may therefore ask whether limited resources would have broader impact in a category such as \texttt{bed\_bath\_table}, which serves more orders. The analysis must jointly inspect service risk, product revenue, and business volume rather than equate one extreme metric with the best intervention target.

The running interaction proceeds as follows:

\begin{enumerate}
\item The analyst asks VerbalVis to restrict the dashboard to RJ and the specified date range.
\item The system constructs coordinated category views for low-score ratio, delivery time, product revenue, and order count using one common product-revenue Top~15 set.
\item The assistant begins explaining a preliminary risk-based interpretation.
\item The analyst interrupts: ``Wait---do not decide from low ratings and delivery time alone. Take revenue and order volume into account.''
\item VerbalVis stops the obsolete audio and cancels the active model response. Late events from that response are ignored.
\item The latest utterance is interpreted against the retained RJ scope and existing views. Any new valid function calls are executed in a closed tool batch, and the updated dashboard state is returned to Qwen before a new explanation is generated.
\end{enumerate}

This example illustrates the distinction between \emph{established state} and an \emph{active response}. The RJ and date filters are established and remain useful. The unfinished recommendation is response-dependent and becomes obsolete. The new request revises the working hypothesis---from ``highest risk should receive priority'' to ``priority must balance risk and impact''---without requiring the user to reconstruct the prior scope.

\section{VerbalVis Design}
\label{sec:design}

VerbalVis is organized around a continuous loop in which speech expresses analytical direction and visualization preserves evidence:

\begin{equation}
\begin{split}
\text{Visual Observation}
&\rightarrow \text{Spoken Analytical Request}\\
&\rightarrow \text{Tool-Supported Analysis}\\
&\rightarrow \text{Dashboard Update}\\
&\rightarrow \text{New Visual Observation}.
\end{split}
\label{eq:interaction-loop}
\end{equation}

\subsection{Context-Grounded Interpretation}
\label{sec:design-grounding}

Requests such as ``only keep RJ,'' ``compare that with order volume,'' or ``highlight week 48'' depend on entities and constraints already visible or recently discussed. VerbalVis therefore combines recent dialogue with compact dashboard metadata: active filters, highlighted values, view identifiers, titles, encodings, and summary statistics. The dashboard functions as part of the conversational context rather than a passive output surface (DR1).

VerbalVis does not maintain a formal symbolic representation of goals or hypotheses. The three revision constructs are used for analysis and design reasoning; runtime reference resolution is delegated to the realtime model using conversation history and the injected dashboard state.

\subsection{Composable Analytical Actions}
\label{sec:design-actions}

An analytical utterance may require several operations. The RJ request, for example, combines a geographic constraint, a time interval, a common Top~15 category set, and four metrics. VerbalVis exposes schema-grounded tools for scope updates, aggregation, group comparison, coordinated category comparison, view creation and modification, deletion, highlighting, inspection, dashboard summarization, and undo (DR2).

Schema grounding constrains the model to dataset fields, supported metrics, and line, bar, or scatter visualizations. The model supplies structured arguments rather than arbitrary SQL or chart code. Server-side handlers normalize and validate these arguments before execution.

\subsection{Immediate Floor Yielding}
\label{sec:design-overlap}

VerbalVis adopts an immediate R-A yielding policy (DR3). When Qwen Semantic VAD reports speech onset during assistant output, the system treats the event as a request for the floor. It stops browser playback and sends \texttt{response.cancel}. The policy is intentionally simple: a short acknowledgement or noise event may also stop the response, and the system does not resume the exact unheard suffix. The completed utterance is subsequently interpreted in context to decide whether it is a repair, clarification, follow-up, or analytical redirection.

\subsection{Response-Level Coordination and Tool Boundary}
\label{sec:design-coordination}

The central coordination unit is a realtime model response. The backend stores a \texttt{current\_response\_id}, and the frontend stores the corresponding active playback response. Audio and transcript deltas are accepted only for that identifier. On barge-in, the identifier is cleared, the response is added to an interrupted set, playback stops, and a cancel event is sent to Qwen. If a later \texttt{response.done} event arrives for the interrupted response, it is discarded.

Function calls are not executed incrementally as soon as partial arguments appear. They are extracted from \texttt{response.done} only when the response is still current and has completed normally. This is the tool-admission check: unfinished calls belonging to a cancelled response never enter the tool layer.

Once an admitted tool batch starts, VerbalVis changes policy. The batch is a non-preemptive, input-closed boundary. Browser audio capture is blocked from transmission, backend audio chunks are rejected, and the provider input buffer is cleared. Tools execute sequentially, with later calls skipped after a blocking failure. Running tools are not cancelled and their mutations are not rolled back. After the batch, each function output contains the current dashboard state; one \texttt{response.create} then requests a spoken explanation grounded in the completed results (DR4).

\subsection{Persistent Visual Workspace}
\label{sec:design-workspace}

Speech is transient and sequential, whereas visual analysis often requires simultaneous comparison. VerbalVis therefore maintains a persistent dashboard with four base views: monthly order trend, review-score distribution, orders by state, and product-revenue Top~15 categories. Additional line, bar, and scatter views can be created and modified during conversation. All views occupy equal responsive grid cells; filters and coordinated comparisons preserve a consistent analytical scope.

Highlighting can focus a view or a data value such as \texttt{2017-W48}, \texttt{office\_furniture}, or their intersection. This makes the spoken explanation inspectable: the user can verify whether the claimed anomaly or comparison is visible in the dashboard.

\subsection{Interaction Feedback}
\label{sec:design-feedback}

The interface exposes connection state, microphone state, active filters, view identifiers, and a compact transcript. User and assistant text occupy one row per utterance or response. Tool calls appear as one-line summaries and expand only to show the tool name and parameters. An interrupted assistant response is retained in the transcript with a small marker, while its audio no longer plays.

\section{System Implementation}
\label{sec:system}

VerbalVis is a browser-based, single-user research prototype comprising a Vue~3 frontend, a FastAPI backend, Qwen~3.5 Omni Plus Realtime, schema-grounded tools, an in-memory DuckDB database, and Vega-Lite visualizations.

\subsection{Session and Audio Pipeline}
\label{sec:audio-pipeline}

One page lifecycle corresponds to one browser WebSocket, one backend session manager, and one Qwen Realtime session. Starting or stopping the microphone changes whether PCM chunks are sent; it does not create or destroy the conversational session.

The browser captures 16-kHz mono PCM16 audio using an \texttt{AudioWorklet}. Chunks are approximately 40~ms, base64 encoded, and transmitted over the application WebSocket. Qwen Semantic VAD detects speech onset and offset. Assistant audio is returned as 24-kHz PCM16 deltas and scheduled through the Web Audio API.

The frontend enforces a hard playback invariant: audio from at most one assistant response may be scheduled at any time. Beginning a new response stops all sources associated with the previous response. Natural completion is reported only after \texttt{response.done} has been received and every scheduled source for that response has finished.

\subsection{Realtime Event Coordination}
\label{sec:realtime-coordination}

The backend forwards audio and transcript deltas only when their response identifier matches \texttt{current\_response\_id}. Speech onset triggers immediate playback stop and \texttt{response.cancel}. The active identifier is cleared before late events can be processed, so old audio cannot be played alongside a newer response.

A completed response may contain zero or more function calls. Calls are normalized, validated, and executed sequentially. The official Qwen tool flow is followed: all \texttt{function\_call\_output} items are inserted first, and one \texttt{response.create} is sent after the batch. During this interval, microphone input remains blocked until the post-tool response starts, preventing user audio from racing with the follow-up response creation.

\subsection{Schema-Grounded Analytical Tools}
\label{sec:system-tools}

The active implementation registers eleven tools (Table~\ref{tab:tools}). The tools operate on shared dashboard state rather than returning unconstrained chart code.

\begin{table*}[t]
\centering
\caption{Model-facing analytical tools in VerbalVis.}
\label{tab:tools}
\small
\begin{tabularx}{\textwidth}{p{0.22\textwidth}p{0.27\textwidth}X}
\toprule
\textbf{Tool} & \textbf{Principal arguments} & \textbf{Purpose} \\
\midrule
\texttt{update\_analysis\_scope} & operation, filters, fields & Replace, add, remove, or clear shared filters. \\
\texttt{aggregate\_data} & group\_by, metrics, filters, sort, limit & Compute grouped evidence without creating a chart. \\
\texttt{compare\_selected\_groups} & dimension, values, metrics, time grain & Compare explicitly named states, categories, scores, or times. \\
\texttt{compare\_category\_metrics} & mode, Top-$N$, ranking, metrics, scope & Create coordinated views for one common category set. \\
\texttt{create\_visual} & chart type, $x$, $y$, series, Top-$N$ & Create one validated line, bar, or scatter view. \\
\texttt{update\_visual} & view id and changed encodings & Modify a view while preserving its identifier. \\
\texttt{delete\_visual} & view id & Remove one dashboard view. \\
\texttt{highlight\_visual} & view ids, value or field--value expression & Direct attention to views or marks without filtering data. \\
\texttt{inspect\_visual} & view id, series, $x$ values, Top-$k$ & Read targeted evidence from a current view. \\
\texttt{summarize\_dashboard} & none & Return filters, views, highlights, and compact statistics. \\
\texttt{undo\_last\_action} & none & Restore the preceding completed dashboard-changing state. \\
\bottomrule
\end{tabularx}
\end{table*}

\subsection{Data and Metric Semantics}
\label{sec:data-semantics}

The Olist data layer contains an order-level table with 96,478 delivered orders and an item-level table with 110,197 order-item records. The order-level table supports temporal, review, state, and delivery analyses; the item-level table supports category and revenue aggregation.

Metric definitions are explicit. Low-score orders satisfy \texttt{review\_score <= 2}. Product revenue is \texttt{SUM(price)} and excludes freight. Category delivery metrics are computed at one row per \texttt{order\_id + product\_category}, preventing orders containing multiple items from the same category from receiving extra weight. Late delivery indicates that the actual delivery date is after the estimated date.

The backend is authoritative for filters, views, highlights, and undo history. After a state-changing tool succeeds, the complete view collection and compact dashboard state are sent to the frontend. This whole-state synchronization is appropriate for the present single-user prototype but does not provide isolated concurrent sessions.

\subsection{Logging and Validation}
\label{sec:logging}

Each connection produces an \texttt{events.jsonl} trace and a \texttt{conversation.jsonl} transcript in a directory that includes the analysis identifier. Logged events include response creation, speech onset, cancellation requests, tool calls, tool results, dashboard state, and timing metadata. The repository also contains deterministic validation scripts for tool behavior, Task~A and Task~B data semantics, response-session lifecycle, in-chart highlighting, responsive layout, and frontend build integrity.

The traces support reconstruction of the implemented mechanism, but they do not constitute a human-performance evaluation. In particular, the present work does not report comparative measures of workload, perceived control, task time, or insight quality.

\section{Case Study: SP Weekly Operational Risk}
\label{sec:case-study}

We use a separate SP scenario as the main case study. The objective is to examine whether a proposed operational-risk period is supported consistently across product categories and metrics. The scope is SP between October 2017 and May 2018. Categories are selected once as the product-revenue Top~5 and reused across four weekly views: order count, low-score ratio, average delivery days, and late-delivery ratio. Week \texttt{2017-W48} is treated as the focal week.

\subsection{Case Procedure}

The analyst begins with the question: ``Is week 48 a common operational peak for our major categories, and does high demand coincide with customer-experience risk?'' VerbalVis atomically applies the SP and date filters and invokes \texttt{compare\_category\_metrics} in \texttt{weekly\_trends} mode. The tool returns four coordinated multi-series views using the same category set, together with each category's focal-week value and peak week for every metric.

During the initial explanation, the analyst interrupts: ``Do not assume the order peak is the risk peak. Compare low ratings and late delivery separately.'' The old spoken response stops. The already established SP scope and coordinated views remain visible. The revised request is interpreted as a working-hypothesis revision: high order volume is no longer treated as sufficient evidence of service risk. VerbalVis then highlights the focal week and inspects the relevant series rather than rebuilding the complete workspace.

The analyst subsequently narrows attention to the category whose demand and risk patterns diverge most strongly. This is a scope refinement, not a new top-level goal. The dashboard retains the common Top~5 category set, allowing the analyst to compare the focal category against its peers without changing the selection basis mid-analysis.

\subsection{Analytical Observations}

The coordinated views expose three useful patterns.

First, \texttt{2017-W48} is not necessarily a simultaneous peak in order count, low-score ratio, delivery time, and late ratio for every category. The visual comparison prevents a demand spike from being interpreted automatically as a customer-experience crisis.

Second, cross-metric rankings can diverge. In the case trace, one category combines the highest product revenue with a comparatively low low-score ratio and a relatively small order count. This pattern suggests that revenue may be driven by higher value per order rather than broad demand. The dashboard makes this interpretation inspectable because revenue, volume, and service metrics remain visible together.

Third, the category with the strongest risk intensity need not be the category where an intervention reaches the largest number of orders. This observation connects the SP case study to the RJ running example: prioritization requires a decision criterion that balances severity and impact rather than selecting the largest value in one chart.

\subsection{Mechanism Observations}

The case trace demonstrates the implemented coordination behavior:

\begin{itemize}
\item When the analyst interrupts, the obsolete assistant audio stops and cannot overlap with the next response.
\item The interrupted response's late audio, transcript, and completion events are discarded by response-identifier checks.
\item Established filters and views remain available after the hypothesis revision.
\item Valid tool calls run in a sequential, input-closed batch; already running tools are not cancelled.
\item The post-tool answer is generated only after function outputs and the latest dashboard state have been returned to Qwen.
\end{itemize}

These observations show that the prototype supports a coherent end-to-end analytical trace. They do not show that full-duplex speech improves human performance relative to another interface.

\section{Discussion}
\label{sec:discussion}

\subsection{Why Combine Voice and Visualization?}

The case study illustrates why the two modalities are complementary. Voice supports rapid expression of changes in analytical direction. An analyst can correct a premise or add a criterion without moving to a text field or waiting for the assistant to finish. Visualization supports simultaneous comparison, persistent memory, and anomaly localization. It allows users to verify whether a spoken summary is consistent with the data and to notice combinations the assistant did not verbalize.

The design implication is not that every analytical interaction should be spoken. Rather, voice is useful for steering, while the dashboard is useful for inspection. VerbalVis treats the visual workspace as part of the conversation state rather than as a disposable response artifact.

\subsection{Response-Level Consistency Rather Than Transactional Tools}

The strongest guarantee in the current implementation is response-level consistency. An interrupted response cannot continue playing and cannot submit function calls that were not admitted before cancellation. This is narrower than transactional tool invalidation. Once a valid tool batch starts, VerbalVis blocks further input and allows the batch to finish. It does not cancel database queries, roll back dashboard mutations, or attach an epoch to every state update.

This boundary simplifies the system and matches the present single-user case-study goal, but it also means that users cannot revise the analysis while a tool batch is running. A future implementation could compute changes against an isolated state snapshot and commit them only after a second validity check, or could expose cancellable query operators.

\subsection{Limitations}

The evidence is case-based. The formative inquiry involved four participants and informed design requirements, but the SP case study does not estimate usability, workload, perceived control, or analytical quality across a population. The running example and case study use one e-commerce dataset, so generalization to clinical, scientific, or financial settings remains open.

The immediate R-A policy may stop an assistant response for a backchannel or false-positive speech event. The prototype does not classify backchannels, repairs, clarifications, and analytical revisions before yielding the floor. Qwen may also retain generated but unheard content in its provider-side conversation because the implementation stops client playback and cancels the response but does not truncate the assistant item at the exact playback position.

The backend stores dashboard state in process-level memory and supports one active analytical session. Tool execution is non-preemptive, dashboard updates are not versioned, and the frontend does not independently reject out-of-order state versions. These constraints are appropriate for a controlled prototype but not for production deployment.

\subsection{Future Work}

Future work should investigate transaction-aware tool execution, explicit dashboard versions, richer interruption interpretation, and isolated multi-user state. A subsequent human evaluation could examine how analysts actually use barge-in and whether the coordination mechanism changes perceived control or analytical behavior, but such claims are outside the evidence reported here.

\section{Conclusion}
\label{sec:conclusion}

We presented VerbalVis, a full-duplex conversational visual analytics system for analytical intent revision. The system combines rapid spoken steering with a persistent visual workspace for comparison and anomaly inspection. Its implemented coordination mechanism uses response-level cancellation, response-identifier filtering, current-response tool admission, a non-preemptive sequential tool boundary, and post-execution dashboard-state reinjection.

A four-participant formative inquiry informed the design. The RJ resource-allocation scenario served as a running example, and the SP weekly operational-risk case study demonstrated how the system supports cross-metric reasoning, working-hypothesis revision, scope refinement, and synchronized speech--tool--dashboard updates. The contribution is a case-based system account rather than a causal user-study result. VerbalVis shows how a realtime speech model can be embedded within a visual analytics controller that keeps spoken output, accepted analytical actions, and the visible workspace aligned at the response level while making its non-transactional tool boundary explicit.

\bibliographystyle{IEEEtran}
\bibliography{main}

\end{document}
